%%% MA_05_Monetary_Policy_Inflation.tex -- Session 5

\section{Banking and Inflation Control}

\begin{frame}{Banking}
    \showmytikz{../shared/tikz/tikz-bankingsector}
    \mycomment{
    The central bank provides central bank money (reserves) to banks, against collateral and at its policy rate. The banks lend to businesses (and households, e.g.\ mortgages), creating deposits, on the condition of interest payments.

    The businesses employ consumers as workers, and produce goods worth more than the credit they have received, paying salaries and interest.

    So the system works only if investment earns more than the interest rate---this links interest rates to growth. The entire system hinges on the central bank setting the policy rate.
    }
\end{frame}

%%%%%%%%%%%%%%%%%%%%%%%%%%%%%%%%%%%%%%%%%%%%%%%%%%%%%%%%%%%%%%%%%%%%%%%%%%%%%%
%%% (F) LIMITS AND REGULATION
%%%%%%%%%%%%%%%%%%%%%%%%%%%%%%%%%%%%%%%%%%%%%%%%%%%%%%%%%%%%%%%%%%%%%%%%%%%%%%

\begin{frame}{What Limits Money Creation?}
    \mytext{
    If not the reserve requirement---what stops banks from lending without limit?
    \begin{itemize}
        \item the \textbf{price}: the central bank's policy rate (costs of refinancing)
        \item \textbf{demand}: there must be creditworthy borrowers who want loans
        \item \textbf{capital requirements} (Basel III): losses must be covered by equity
        \item \textbf{liquidity rules}: deposits can leave the bank
        \item \textbf{supervision}
    \end{itemize}
    }
    \mycomment{
    So the central bank controls money creation mainly through the interest rate, not through the quantity of money. This is why the textbook multiplier is a poor description of modern banking \cite{mcleay2014money}.
    }
\end{frame}



\begin{frame}{Reserve Requirements vs.\ Capital Requirements}
    \mytable{
        \renewcommand{\arraystretch}{1.2}
        \begin{tabularx}{\textwidth}{>{\raggedright\arraybackslash}p{2.8cm}
                                     >{\raggedright\arraybackslash}X
                                     >{\raggedright\arraybackslash}X}
        \toprule
         & \textbf{Reserve requirement} \newline (Mindestreserve) & \textbf{Capital requirement} \newline (Eigenkapitalanforderung) \\
        \midrule
        \textbf{bank must hold} & central bank money (reserves) & own funds (equity) \\
        \textbf{measured against} & deposits & loans, weighted by risk \\
        \textbf{size} & 1\,\% (euro area) & ca.\ 10\,\% \\
        \textbf{protects against} & lack of liquidity & losses \\
        \textbf{binding today?} & no & yes \\
        \bottomrule
        \end{tabularx}
 
        \bigskip
 
        {\footnotesize Careful: central bank reserves (Zentralbankguthaben) are an \textit{asset}; \enquote{Rücklagen} are part of \textit{equity}.}
    }
\end{frame}
 
\begin{frame}{Capital Requirements}
    \mytext{
    Banks must finance part of their loans with their own money (equity):
    \begin{equation}
        \colorbox{yellow!40}{$\dfrac{\text{equity}}{\text{risk-weighted loans}} \geq \text{ca.\ } 10\,\%$}
    \end{equation}
    \begin{itemize}
        \item riskier loans count more: business loans fully, mortgages less, government bonds not at all
        \item example: with 10 of equity, a bank can hold about 100 of business loans
        \item to lend more, the bank needs more equity (retain profits, issue shares)
        \item losses are borne by the bank's owners first, not by depositors or taxpayers
    \end{itemize}
    }
    \mycomment{
    The international standard is set by the Basel Committee (Basel III, agreed after the financial crisis) and is EU law. The exact requirement differs by bank: a minimum of 8\,\% plus several buffers.
 
    This is the real limit on bank lending today---not the reserve requirement.
    }
\end{frame}


\begin{frame}{Control Money Supply}
    \mycomment{
        Since banks can create money out of nothing, there is a temptation to do so. But that undermines the stability of the currency through inflation, and undermines trust.

        Thus, banks are strictly regulated and supervised. Only the central bank may issue bank notes. Private banks participate in the creation of money, but under oversight.

        Central banks are made independent so that price stability is not sacrificed to short-term political goals (e.g.\ financing the government).
    }
    \mytext{
        \begin{itemize}
            \item How to stop banks from creating too much money?
            \item How to control the value of money?
            \item How to control inflation?
        \end{itemize}
    }
\end{frame}

\begin{frame}{Control Money Supply}
    \mytext{
    Banks are under public supervision:
    \begin{itemize}
        \item Euro area: ECB (Single Supervisory Mechanism, since 2014) for significant banks
        \item Germany: BaFin and Bundesbank for the others
    \end{itemize}

    Tools:
    \begin{itemize}
        \item capital requirements (Basel III)
        \item liquidity requirements
        \item minimum reserve requirements
        \item refinancing through the central bank (at the policy rate)
        \item deposit insurance, resolution rules
    \end{itemize}
    }
\end{frame}

%%%%%%%%%%%%%%%%%%%%%%%%%%%%%%%%%%%%%%%%%%%%%%%%%%%%%%%%%%%%%%%%%%%%%%%%%%%%%%
%%% (G) INTEREST RATES AND MONEY DEMAND
%%%%%%%%%%%%%%%%%%%%%%%%%%%%%%%%%%%%%%%%%%%%%%%%%%%%%%%%%%%%%%%%%%%%%%%%%%%%%%

\begin{frame}{Nominal vs.\ Real Interest Rate}
    \mytext{
    \begin{itemize}
        \item Growth rate of GDP $g=\frac{dY/dt}{Y}$
        \item Inflation $\pi=\frac{dP/dt}{P}$
        \item Nominal interest rate $i$ 
    \end{itemize}    

    Real interest rate (Fisher equation):
    \begin{equation}
        r=\frac{1+i}{1+\pi}-1
    \end{equation}
    \begin{equation}
        \colorbox{yellow!40}{$r \approx i -\pi$ }
    \end{equation}
    }
    \mycomment{
    The approximation is good for small rates. Example: $i=5\,\%$, $\pi=3\,\%$: exact $r=1.94\,\%$, approximately $2\,\%$.

    When lending or investing, what matters is \textit{expected} inflation $\pi^e$ (ex-ante real rate $r = i - \pi^e$). Afterwards, we can compute the realised (ex-post) real rate with actual inflation.

    If capital can move freely, lending money and investing in real capital should earn similar (risk-adjusted) returns. In practice, the link between interest rates and the profit rate (and by extension, economic growth) is weak.
    }
\end{frame}

\begin{frame}{Money Demand}
    \showmytikz[time]{../shared/tikz/tikz-fourfunctionsofmoney}
    \mycomment{
    Why do people hold money (liquidity preference, $L$)?
    \begin{itemize}
        \item to make transactions: more output $Y$ needs more money (compare $MV=PY$)
        \item as precaution and store of value over time
    \end{itemize}
    But holding money has an opportunity cost: the interest foregone on bonds or savings. So money demand \textit{falls} as the interest rate $i$ rises:
    \begin{equation*}
        L = L(\underset{+}{Y},\underset{-}{i})
    \end{equation*}
    We will need this for the $LM$ curve (Session 6).
    }
\end{frame}

\begin{frame}{Money Supply and Money Demand}
\mytext{
    \begin{itemize}
        \item demand for money is demand for credit
        \item[$\rightarrow$] interest rate is the result of the credit market
        \item or: demand for money is demand for liquidity for economic activity
        \item[$\rightarrow$] interest rate is the result of the money market
        \item disputed...
    \end{itemize}
}
\end{frame}

\begin{frame}{Nature of Interest Rates}
\mycomment{
While it is clear that the interest rate is the price of borrowing money, it is less clear where this interest rate actually \enquote{comes from}. Is it set by the monopolist (that is, the issuer of currency, the central bank), or is it a market outcome?

The answer is: both. The central bank sets \textit{short-term} rates, acting like a monopolist for central bank money. But it reacts to the economy, and \textit{long-term} rates are formed on bond markets: they reflect expected future short-term rates plus a premium for risk and maturity.

And to which end does the central bank set the rate---price stability, employment, financing the state?
}
\mytext{
        Where do interest rates come from?
        \begin{itemize}
            \item short-term: set by the central bank
            \item long-term: determined by markets?
            \item to which end?
            \item determined by the growth rate?
        \end{itemize}
}
\end{frame}

%% [moved from MA_4_Money.tex, frame 37]
\begin{frame}{Nature of Interest Rates}
\mytext{
    {Set by the central bank?}
    \begin{itemize}
        \item national central bank / Eurosystem (ECB)
        \item leading central bank (Federal Open Market Committee, US) influences all others
        \item IMF, BIS: no rate setting, but coordination and emergency lending (\enquote{lender of last resort} for states)
    \end{itemize}
}
\end{frame}

\begin{frame}{Nature of Interest Rates}
\mytext{
    {Determined by the growth rate?}
    \begin{itemize}
        \item What determines the growth rate $g$?
        \item What is in the basket of $P$?
        \item Technological progress?
        \item Population growth?
    \end{itemize}

    Benchmark (\enquote{golden rule}): real interest rate $r$ = real growth rate $g$.

    In reality: $r<g$ in the 2010s, $r>g$ in other periods.
    % (Session 1, debt sustainability)
}    
    \mycomment{
    It is not exactly clear what causes the economy to grow. The most obvious is population growth: if output per person is stable, more persons will produce more. On top, technological progress, increased efficiency, and new energy sources lead to increased output per head (Session 11).

    Whether the interest rate reacts to growth or drives growth is a matter of debate.
    }
\end{frame}

%%%%%%%%%%%%%%%%%%%%%%%%%%%%%%%%%%%%%%%%%%%%%%%%%%%%%%%%%%%%%%%%%%%%%%%%%%%%%%
%%% (A) INFLATION: TYPES, CAUSES, COSTS
%%%%%%%%%%%%%%%%%%%%%%%%%%%%%%%%%%%%%%%%%%%%%%%%%%%%%%%%%%%%%%%%%%%%%%%%%%%%%%

\begin{frame}{Types of Inflation}
    \mytable{
        \begin{tabularx}{\textwidth}{>{\raggedright\arraybackslash}p{3.2cm}
                                     >{\raggedright\arraybackslash}X
                                     >{\raggedright\arraybackslash}X}
        \toprule
         & \textbf{Definition: Rising...} & \textbf{Primary measurement} \\
        \midrule
        \textbf{Consumer price inflation} & prices of consumer goods and services & HICP (euro area), CPI/VPI (national) \\
        \addlinespace
        \textbf{Producer price inflation} & prices of goods at the factory gate & PPI \\
        \addlinespace
        \textbf{GDP deflator} & prices of all domestic output & nominal / real GDP \\
        \addlinespace
        \textbf{Asset price inflation} & prices of stocks, houses, and other assets & stock indices (DAX, MSCI World), house price indices \\
        \bottomrule
        \end{tabularx}

        % \bigskip

        % {\footnotesize Without qualification, \enquote{inflation} means a sustained rise of consumer prices; the ECB's target refers to the HICP (Session 2). Some authors (older usage, Austrian school) call money growth itself \enquote{monetary inflation}---today it is treated as a possible \textit{cause} of inflation.}
    }
\end{frame}

\begin{frame}{Demand-Pull Inflation}
    \mycomment{
    Prices can rise because aggregate demand increases---e.g.\ through expansionary fiscal policy, a credit boom, an export boom, or households spending accumulated savings (as after the COVID lockdowns). Fast money growth works through this channel, too.

    Here, we preview a model with an \textit{aggregate} demand curve $AD$, and an \textit{aggregate} supply curve $AS$. We will discuss in much more detail what these actually are, and what influences them (Session 8).

    An increase in aggregate demand leads to higher prices \textit{and} higher output. This is the typical pattern in a boom: inflation together with low unemployment (Phillips curve, Session 7).
    }
    \showmytikz[as,ad,eq,adshift]{../shared/tikz/tikz-asad-inflation}
\end{frame}

\begin{frame}{Cost-Push Inflation}
    \mycomment{
    The overall price level can also increase because of a negative supply shock---for example, less or more expensive energy (oil crises 1973 and 1979, gas prices 2022). In that case, inflation comes together with a \textit{decrease} in output: \textbf{stagflation}.

    This is the harder case for central banks: fighting inflation with higher interest rates deepens the recession.
    }
    \showmytikz[as,ad,eq,asshift]{../shared/tikz/tikz-asad-inflation}
\end{frame}

\begin{frame}{Built-in Inflation}
    \mycomment{
    When wages and other contracts are indexed to past inflation, or when everybody expects inflation and sets wages and prices accordingly, inflation perpetuates itself: a \textbf{wage-price spiral}.

    This is why central banks care so much about inflation \textit{expectations}.
    }
    \showmytikz[]{../shared/tikz/builtininflation}
\end{frame}

\begin{frame}{Inflation, Wages, and Employment}
    \mycomment{
    Inflation is a sustained rise of the general price level. Wages are typically negotiated at longer intervals than prices are set. So at first, inflation lowers real wages, which makes production cheaper. What happens next depends on the labour market.
    }
    \mytext{
    \begin{itemize}
        \item rising prices lower real wages
        \item under full employment:
        \begin{itemize}
            \item unions demand compensation
            \item wages go up
            \item costs and demand go up
            \item prices go up
            \item feedback loop (wage-price spiral)
        \end{itemize}
        \item no full employment:
        \begin{itemize}
            \item more production
            \item more employment
            \item higher income
        \end{itemize}
    \end{itemize}    
    }
\end{frame}

\begin{frame}{Costs of Inflation}
    \mylonglist{
        \begin{itemize}
            \item erosion of purchasing power of money holdings and non-indexed incomes
            \item even when inflation is \textbf{expected}:
            \begin{itemize}
                \item shoe-leather costs: holding less cash, going to  bank more often
                \item menu costs: updating price lists
                \item tax distortion (bracket creep, taxes on nominal gains)
                \item confusion: relative prices harder to read
            \end{itemize}
            \item when inflation is \textbf{unexpected}:
            \begin{itemize}
                \item redistribution of wealth from lenders to borrowers
                \item fixed-income squeeze (non-indexed pensions, wages, rents)
                \item investment uncertainty, higher risk premia, higher borrowing costs
                \item instability
            \end{itemize}
        \end{itemize}
    }
\end{frame}

\begin{frame}{Inflation and Fiscal Drag}
    \mytext{
    Tax rule: 20\,\% on income up to 10,000 EUR, 50\,\% on every euro above.

    \textbf{Before:} income 10,000 EUR $\rightarrow$ tax 2,000 (20\,\%), net 8,000.

    Ten years pass, the tax law is not changed. Prices rise by 50\,\%, nominal wages also by 50\,\%. There is no real growth.

    \textbf{After:} income 15,000 EUR $\rightarrow$ tax $2,000 + 2,500 = 4,500$ (30\,\%), net 10,500.

    In prices of the base year: $10,500/1.5 = 7,000 < 8,000$.

    $\Rightarrow$ Real net income fell by 12.5\,\%, although real gross income did not change.
    }
    \mycomment{
    This is called \enquote{Kalte Progression} in Germany; in English, \enquote{bracket creep} or \enquote{fiscal drag}: inflation pushes taxpayers into higher brackets although they are not better off. Tax revenue (and thus the state quota) rises without any change in the law.

    Real income growth also moves people into higher brackets---but that is the intended effect of a progressive tax. In Germany, the federal government reports on the effect every two years, and the tariff has been adjusted repeatedly to offset it.
    }
\end{frame}

%%%%%%%%%%%%%%%%%%%%%%%%%%%%%%%%%%%%%%%%%%%%%%%%%%%%%%%%%%%%%%%%%%%%%%%%%%%%%%
%%% (B) MONETARY POLICY
%%%%%%%%%%%%%%%%%%%%%%%%%%%%%%%%%%%%%%%%%%%%%%%%%%%%%%%%%%%%%%%%%%%%%%%%%%%%%%

\begin{frame}{Central Bank Mandate: ECB}
    \mycomment{
    In many countries the central bank is independent of the government. The historical background is the experience of governments using monetary policy excessively for their funding needs, causing inflation, currency devaluation and general imbalances.

    The ECB has a single primary objective, in line with Tinbergen's rule. Not all central banks do: the US Fed has a \textbf{dual mandate} (maximum employment and stable prices; target 2\,\% PCE inflation). The Bank of England's target is set by the government; the bank decides only how to reach it.
    }
    \mytext{
    \textbf{Primary mandate} (Art.\ 127 TFEU): maintain price stability. The ECB defines this as (symmetric, over the medium term):
    \begin{equation}
        \pi^*_{\text{HICP}} = 2\,\%
    \end{equation}
    
    \textbf{Secondary mandate:} without prejudice to price stability, support the general economic policies in the EU

    \textbf{Additional responsibilities:} financial stability, banking supervision (since 2014), climate and broader considerations
    }
\end{frame}

\begin{frame}{Instruments of the ECB}
    \mytext{
    \begin{itemize}
        \item \textbf{key interest rates}:
        \begin{itemize}
            \item deposit facility rate (DFR): what banks earn on reserves
            \item main refinancing operations (MRO): what banks pay for weekly loans
            \item marginal lending facility (MLF): overnight emergency loans
        \end{itemize}
        \item \textbf{open market operations}: refinancing operations against collateral, asset purchases
        \item \textbf{minimum reserves}: 1\,\% of certain deposits
        \item \textbf{forward guidance}: communicating future policy
    \end{itemize}
    }
    \mycomment{
    Since the large asset purchases, banks hold far more reserves than required. Money market rates therefore settle near the deposit facility rate, which the ECB uses as its steering rate.
    }
\end{frame}

\begin{frame}{Open Market Operations}
     \mytext{
     Central banks implement their interest rate by buying and selling (or lending against) securities: \textbf{open market operations (OMO).} Traditionally, they supply just enough reserves to keep the money market rate at the target.

     When the policy rate is already near zero, central banks can buy large amounts of bonds outright: \textbf{quantitative easing (QE).} This lowers long-term yields and signals that rates will stay low.

     In the euro area, the central bank may buy government bonds only on the secondary market---direct financing of governments is prohibited (Art.\ 123 TFEU).
     }
    \mycomment{
     Quantitative Easing (QE) is a pre-commitment to buy large, fixed volumes of longer-term assets over an extended period. Central banks used this policy extensively from the 2008 crisis on, and again in the COVID crisis.

     The reserves so created are still in the system: central banks now pay interest on them, which made them post losses when rates rose in 2022--23; reducing the holdings (quantitative tightening) takes years.
    }
\end{frame}

\begin{frame}{Monetary Policy Transmission Mechanism}
    \showmytikz{../shared/tikz/tikz-transmissionmechanism}
    \mycomment{
    The actual transmission mechanism from the official interest rate of the central bank to price developments is complex \cite{ecb_transmission}. 

    In addition, external shocks outside the control of the central bank can occur:

    \begin{itemize}
        \item changes in risk premia
        \item changes in bank capital
        \item changes in the global economy
        \item changes in fiscal policy
        \item changes in commodity prices
    \end{itemize}
    }
\end{frame}

\begin{frame}{Taylor Rule}
    \mytext{
    John \noun{Taylor} (1993) showed that the Fed's behaviour can be described---and, he argued, should be guided---by a simple rule \cite{taylor1993discretion}. The policy rate responds to divergences of
    \begin{itemize}
        \item \textbf{inflation gap: }actual inflation vs.\ inflation target 
        \item \textbf{output gap:} actual GDP vs.\ potential GDP 
    \end{itemize}
    }
    \mycomment{
    Two gaps, one instrument---does this violate Tinbergen's rule? Not necessarily: the output gap matters because it predicts future inflation. But when the two gaps point in different directions (stagflation), the central bank faces a real trade-off.
    }
\end{frame}

\begin{frame}{Taylor Rule}
    \mytext{
    So the rule becomes:
    \begin{equation}
        i_t =  \quad \pi_t+r^* 
        \quad + a_\pi (\pi_t-\pi^*) 
        \quad + a_y\cdot 100\frac{(Y_t-\overline{Y_t})}{\overline{Y_t}}
    \end{equation}
    \begin{equation}
        \begin{split}
            \text{policy rate} \quad  = \quad & \text{inflation} \\
            +\quad & \text{equilibrium real rate} \\
            +\quad & a_\pi \times \text{inflation gap} \\
            +\quad & a_y \times \text{output gap}             
        \end{split}
    \end{equation}
    }
    \mycomment{
    Equivalently, moving $\pi_t$ to the left: the \textit{real} rate $i_t-\pi_t$ should exceed its equilibrium value when inflation is above target or output above potential.

    Taylor's original values: $a_\pi = a_y = 0.5$, $r^*=2\,\%$, $\pi^*=2\,\%$.

    Taylor principle: if inflation rises by one point, the nominal rate rises by $1+a_\pi = 1.5$ points, so the real rate rises, too.
    }
\end{frame}

\begin{frame}{Time Inconsistency Problem}
    \mytext{
    Expectations (of inflation) play a major role in monetary policy, and thus credibility.

    If the central bank announces low inflation and people build their expectations accordingly, the central bank is tempted to deviate from its announcement (e.g.\ to boost employment in the short run).

    People anticipate this. In the long term, it erodes the credibility of the central bank---and inflation ends up higher, without any gain in employment.

    $\Rightarrow$ one argument for independent central banks with a clear mandate.
    }
\end{frame}